% Reporte de la actividad 2A — Motor de búsqueda léxico y BM25
% Compilación: latexmk -pdf reporte.tex   (o: pdflatex reporte.tex, dos veces)
\documentclass[10pt,a4paper]{article}

\usepackage[T1]{fontenc}
\usepackage[utf8]{inputenc}
\usepackage[spanish,es-nodecimaldot]{babel}
\usepackage[margin=1.9cm,top=1.7cm,bottom=1.7cm]{geometry}
\usepackage{amsmath}
\usepackage{booktabs}
\usepackage{array}
\usepackage{xcolor}
\usepackage{microtype}
\usepackage{enumitem}
\usepackage[hidelinks]{hyperref}

\definecolor{acento}{HTML}{1F6FEB}
\newcommand{\term}[1]{\textcolor{acento}{\texttt{#1}}}

\setlist[itemize]{topsep=2pt,itemsep=1pt,parsep=0pt,leftmargin=1.2em}
\setlist[enumerate]{topsep=2pt,itemsep=1pt,parsep=0pt,leftmargin=1.5em}
\setlength{\parskip}{2pt}
\renewcommand{\arraystretch}{1.05}
\usepackage{titlesec}
\titlespacing*{\section}{0pt}{6pt}{3pt}
\titlespacing*{\paragraph}{0pt}{4pt}{0.6em}

\title{\vspace{-1.4cm}\textbf{Motor de búsqueda léxico y BM25}\\[2pt]
\large Análisis de aciertos y del límite de la búsqueda exacta}
\author{David Segundo García}
\date{Actividad 2A \textperiodcentered\ \today}

\begin{document}
\maketitle
\vspace{-0.7cm}

\section{Sistema evaluado}

Se implementó desde cero, en un \emph{notebook} de Jupyter, un buscador léxico completo sobre un
corpus propio de \textbf{120 documentos en español} repartidos en \textbf{12 temáticas}
(inteligencia artificial, astronomía, salud, gastronomía, deportes, economía, medio ambiente,
historia, arte, transporte, ciberseguridad y viajes). El sistema tiene tres piezas:

\begin{enumerate}
  \item \textbf{Preprocesamiento.} Minúsculas $\rightarrow$ tokenización con expresión regular que
        descarta puntuación $\rightarrow$ filtro de tokens de un carácter $\rightarrow$ eliminación
        de 368 \emph{stopwords} $\rightarrow$ \emph{stemming} Snowball para español $\rightarrow$
        supresión de tildes en la raíz. La \textbf{misma} función se aplica al corpus y a la
        consulta, condición sin la cual los términos no serían comparables.
  \item \textbf{Índice invertido} con diccionarios nativos: \term{término $\rightarrow$ \{documento: tf\}},
        más las longitudes $|D|$, la frecuencia documental $\mathit{df}$ y la longitud promedio
        \emph{avgdl}.
  \item \textbf{Ranking BM25} con $k_1 = 1.5$, $b = 0.75$ e IDF de Robertson en su variante de
        Lucene, verificado contra el cálculo manual de la fórmula.
\end{enumerate}

\[
\text{score}(D,Q)=\sum_{i=1}^{n}\text{IDF}(q_i)\cdot
\frac{f(q_i,D)\,(k_1+1)}{f(q_i,D)+k_1\!\left(1-b+b\dfrac{|D|}{\text{avgdl}}\right)},
\qquad
\text{IDF}(q_i)=\ln\!\left(\frac{N-n(q_i)+0.5}{n(q_i)+0.5}+1\right)
\]

\begin{table}[h]
\centering\small
\begin{tabular}{lr@{\hspace{2.2em}}lr}
\toprule
\textbf{Corpus} & & \textbf{Índice} & \\
\midrule
Documentos            & 120   & Términos únicos (vocabulario) & 1\,940 \\
Temáticas             & 12    & Entradas de \emph{postings}   & 3\,608 \\
Palabras crudas       & 6\,369 & Longitud promedio \emph{avgdl} & 32.50 \\
Tokens indexables     & 3\,900 & Rango de $|D|$                & 24 -- 45 \\
Reducción tras limpiar & 38.8\,\% & Términos con $\mathit{df}=1$ & 1\,211 (62.4\,\%) \\
\bottomrule
\end{tabular}
\end{table}

\section{Metodología de evaluación}

Se definieron 14 consultas con su \textbf{categoría esperada} como criterio de relevancia. Sobre
ellas se calcularon tres métricas: \textbf{P@5} (estricta, con denominador 5), \textbf{P@k efectiva}
(sobre los documentos realmente recuperados) y \textbf{MRR}.

\begin{table}[h]
\centering\small
\begin{tabular}{lccc}
\toprule
\textbf{Métrica} & \textbf{Valor} & \textbf{Interpretación} \\
\midrule
MRR                    & \textbf{1.000} & El documento correcto quedó primero en las 14 consultas \\
P@k efectiva (promedio) & 0.682 & Dos de cada tres resultados devueltos son de la categoría esperada \\
P@5 estricta (promedio) & 0.500 & Penaliza los huecos: varias consultas recuperan menos de 5 documentos \\
\bottomrule
\end{tabular}
\end{table}

La brecha entre ambas precisiones no indica error de ranking sino \textbf{cobertura}: con solo diez
documentos por temática, consultas muy específicas como \emph{«nixtamalización del maíz para
tortillas»} recuperan un único documento —el correcto— y la P@5 estricta las castiga con 0.2. La
categoría es además un criterio duro: en \emph{«energía solar y eólica renovable»} los documentos
de transporte sobre baterías y motores eléctricos cuentan como fallo aunque sean pertinentes. Las
cifras son, por tanto, una cota inferior de la calidad real.

\section{Casos en los que el buscador fue preciso}

\subsection*{Caso 1: términos raros — «nixtamalización del maíz para tortillas»}

Las tres raíces de la consulta (\term{nixtamaliz}, \term{maiz}, \term{tortill}) tienen
$\mathit{df}=1$, es decir, aparecen en un único documento del corpus. Su IDF es el máximo posible
(4.39) y ningún otro documento puede competir:

\begin{table}[h]
\centering\small
\begin{tabular}{lrrrr}
\toprule
\textbf{Término} & \textbf{tf} & \textbf{df} & \textbf{IDF} & \textbf{Aporte BM25} \\
\midrule
\term{nixtamaliz} & 2 & 1 & 4.3903 & 6.1802 \\
\term{maiz}       & 2 & 1 & 4.3903 & 6.1802 \\
\term{tortill}    & 1 & 1 & 4.3903 & 4.3010 \\
\midrule
\multicolumn{4}{r}{\textbf{Total}} & \textbf{16.6614} \\
\bottomrule
\end{tabular}
\end{table}

El desglose exhibe además la \textbf{saturación} de $k_1$: duplicar la frecuencia (de $tf=1$ a
$tf=2$) sube el aporte de 4.30 a 6.18, no a 8.60. El décimo hallazgo de una palabra no vale diez
veces el primero, que es exactamente lo que distingue a BM25 de un TF-IDF ingenuo.

\subsection*{Caso 2: acumulación de evidencia — «ransomware respaldo y cifrado»}

Ningún término es tan raro como en el caso anterior, pero los tres coinciden en el mismo documento
(\emph{Ransomware y recuperación}, score 14.04) mientras que el resto de candidatos acierta solo
uno o dos. La suma de aportes separa con claridad al primero: BM25 premia la
\textbf{cobertura de la consulta}, no la repetición de un único término.

\subsection*{Caso 3: normalización por longitud — «energía»}

Con un solo término en la consulta, dos documentos tienen $tf=2$: \emph{Energía solar fotovoltaica}
($|D|=30$) y \emph{Materia oscura y energía oscura} ($|D|=37$). Con $b=0.75$ gana el documento más
corto (4.283 frente a 3.999), porque concentra la misma evidencia en menos texto. Al fijar $b=0$
ambos empatan en 4.177: la normalización por \emph{avgdl} es lo único que los separa.

\section{Caso de fallo por búsqueda exacta: los sinónimos}

El corpus se redactó deliberadamente usando \emph{automóvil}, \emph{médico} y \emph{computadora}, y
evitando por completo \emph{coche}, \emph{doctor} y \emph{ordenador}. BM25 es un modelo
\textbf{léxico}: puntúa coincidencia de cadenas, no de significados. Un término ausente del
vocabulario tiene $\mathit{df}=0$, su IDF se define como 0 y su aportación al ranking es
exactamente cero. El resultado son \textbf{tres modos de fallo distintos}, de gravedad creciente:

\begin{table}[h]
\centering\small
\begin{tabular}{@{}p{3.1cm}p{2.0cm}rp{5.1cm}r@{}}
\toprule
\textbf{Consulta} & \textbf{Raíces} & \textbf{df} & \textbf{Primer resultado} & \textbf{Score} \\
\midrule
coche eléctrico     & coch, electr & 0, 5 & Automóviles eléctricos y autonomía & 5.017 \\
\emph{automóvil eléctrico} & \emph{automovil, electr} & \emph{2, 5} & \emph{Automóviles eléctricos y autonomía} & \emph{10.373} \\
\midrule
doctor              & doctor & 0 & \textcolor{red}{--- sin resultados ---} & 0.000 \\
\emph{médico}       & \emph{medic} & \emph{1} & \emph{Prevención de la diabetes tipo 2} & \emph{4.187} \\
\midrule
ordenador           & orden & 2 & \textcolor{red}{Chocolate desde el cacao} & 4.077 \\
\emph{computadora}  & \emph{comput} & \emph{4} & \emph{Visión por computadora} & \emph{4.460} \\
\bottomrule
\end{tabular}
\end{table}

\paragraph{1. Degradación silenciosa.} \emph{«coche eléctrico»} sí devuelve resultados, pero por el
motivo equivocado: \term{coch} aporta 0.0 y todo el orden lo decide \term{electr}. El score del
documento correcto cae de 10.373 a 5.017, y compite en igualdad con documentos sobre energía solar
o baterías que también contienen \emph{eléctrico}. Basta un corpus algo mayor para que el documento
buscado deje de aparecer en el Top-5.

\paragraph{2. Cero resultados.} La consulta \emph{«doctor»} devuelve una lista vacía: ningún término
aporta puntuación y el conjunto de candidatos queda vacío. El corpus habla de médicos, de
telemedicina y de especialistas en varios documentos, pero el buscador no tiene forma de saberlo.

\paragraph{3. Falso positivo por colisión de raíces (el más grave).} \emph{«ordenador»} es reducido
por el \emph{stemmer} a \term{orden}, raíz que el corpus usa en «el templado \emph{ordena} los
cristales de manteca de cacao». El primer resultado es un documento sobre chocolate, presentado con
la misma confianza que cualquier acierto. Lo mismo ocurre con \emph{«cita con el doctor»}: la raíz
\term{cit} aparece en «las cuatro grandes \emph{citas} del calendario» y el buscador responde con un
documento de tenis (score 3.801). El \emph{stemming}, que en general mejora el recall, aquí fabrica
una coincidencia que no existe. \textbf{Un resultado equivocado es peor que ningún resultado},
porque no hay señal alguna de que la consulta se entendió mal.

\section{Mitigaciones evaluadas}

\paragraph{Expansión de consulta con diccionario de sinónimos (implementada).} Antes de puntuar se
añaden a la consulta las variantes conocidas de cada término. No requiere reconstruir el índice y su
costo es prácticamente nulo. Recupera los tres casos: \emph{coche eléctrico} pasa de 5.017 a
\textbf{10.373}, \emph{cita con el doctor} deja de responder tenis para devolver un documento
médico (4.187) y \emph{ordenador} recupera \emph{Visión por computadora} (4.460). Su límite es
estructural: solo cubre lo que alguien anticipó; \emph{«vehículo de cero emisiones»} seguiría
fallando.

\paragraph{Recuperación densa e híbrida (discutida).} Representar consulta y documentos como
vectores densos sitúa \emph{coche} y \emph{automóvil} próximos en el espacio semántico aunque no
compartan una sola letra.

\begin{table}[h]
\centering\small
\begin{tabular}{@{}lcccc@{}}
\toprule
\textbf{Enfoque} & \textbf{Sinónimos} & \textbf{Términos raros} & \textbf{Costo} & \textbf{Explicabilidad} \\
\midrule
BM25 (léxico)            & No & Excelente & Muy bajo, CPU & Total \\
Denso (\emph{embeddings}) & Sí & Débil     & Alto          & Baja \\
Híbrido (BM25 + denso)   & Sí & Buena     & Medio-alto    & Media \\
\bottomrule
\end{tabular}
\end{table}

\section{Conclusiones}

\begin{enumerate}
  \item El preprocesamiento fija el techo del sistema: redujo 6\,369 palabras a 3\,900 términos
        indexables y es lo que permite que \emph{baterías} recupere documentos que escriben
        \emph{batería}. Cualquier asimetría entre el camino del corpus y el de la consulta rompe la
        recuperación sin producir ningún error visible.
  \item BM25 acierta de forma consistente cuando consulta y documento comparten vocabulario:
        MRR de 1.000 en las 14 consultas evaluadas. Los aciertos más nítidos son los de términos con
        $\mathit{df}$ bajo, donde el IDF domina la suma de aportes.
  \item Los tres componentes de la fórmula son verificables por separado, pero su efecto agregado en
        un corpus de 120 documentos de longitud homogénea es pequeño: conviene no sobreinterpretar
        una tabla de sensibilidad de $k_1$ y $b$ calculada a esta escala.
  \item El fallo por sinónimos no es un defecto de implementación sino el límite del modelo, y su
        forma más peligrosa no es la lista vacía sino el falso positivo por colisión de raíces.
  \item La elección de mitigación es una decisión de costo: un diccionario de sinónimos resuelve lo
        previsible con costo nulo; la recuperación densa resuelve el caso general a cambio de un
        modelo, un índice vectorial y una pérdida de explicabilidad. La práctica actual —y la razón
        por la que BM25 sigue vigente treinta años después— es combinar ambas señales en un
        recuperador híbrido.
\end{enumerate}

\end{document}
